\documentclass{exam}
\usepackage{../../../mypackages}
\usepackage{../../../macros}

\title{Interrogation N°3}
\author{N. Bancel}
\date{4 Décembre 2025}

% Mise en forme des solutions en bleu
\SolutionEmphasis{\color{blue}}
\renewcommand{\solutiontitle}{\noindent}

\begin{document}

\textbf{Collège Lycée Suger}
\hfill
\textbf{Mathématiques} \\

\textbf{Année 2025-2026}
\hfill
\textbf{Classe de 6ème} \\

{\let\newpage\relax\maketitle}

\begin{center}
\textbf{\textcolor{red}{Durée : 1h15}} \\
\textbf{\textcolor{red}{Chaque question qui commence par [Justifier] attend une rédaction : \\
Pour ce type de question, une réponse sans phrase de rédaction rapporte 0 point.}} \\
Une copie mal tenue sera pénalisée.

\end{center}

\section*{Exercice 1 - Cours (7 points)}

\begin{questions}
  \question[4] Médiatrice 
    \begin{parts}
      \part[1.5] Donner la définition précise de la médiatrice d’un segment. Un vocabulaire mathématique est attendu
      \part[0.5] Quelle est la propriété principale des points d'une médiatrice ?
      \part[2] Avec deux méthodes différentes, tracer la médiatrice d'un segment $AB$ de longueur $6$ cm. Inclure toutes les légendes nécessaires. \textbf{Il est attendu de nommer la méthode. Mais pas besoin de détailler les étapes}
    \end{parts}
  \question[1.5] Triangle : À l'aide de la règle et du compas, construire, en précisant les légendes, le triangle $ABC$ tel que :
\[
AB = 7~\text{cm}, \qquad BC = 10~\text{cm}, \qquad AC = 8.5~\text{cm}.
\]

\question[1.5] Recopier les phrases suivantes sur votre cahier, en complétant les trous.

\begin{compactitem}
    \item Les \quad \ldots \ldots \ldots \quad \ldots \quad et \quad \ldots \quad sont \quad \ldots \quad en $F$.
    \item Le point $F$ est \quad \ldots \ldots des droites $(d)$ et $(d')$.
    \item $F$ \quad \ldots \ldots \quad $(d)$ et $F$ \quad \ldots \ldots \quad $(d')$.
    \item Le \quad \ldots \ldots \quad $[DF]$ a pour \quad \ldots \ldots \quad les points $D$ et $F$
\end{compactitem}

\begin{figure}[H]
  \centering
  \includegraphics[width=0.5\linewidth]{figure_1.jpg}
  \captionsetup{labelformat=empty}
  \caption{\label{} Figure 1}
\end{figure} 
\end{questions}




\section*{Exercice 2 - Tracés géométriques (5 points)}

\begin{questions}
\question[2] Tracer une figure qui respecte les instructions suivantes 

\begin{compactitem}
    \item Tracer une droite $(AB)$.
    \item Placer un point $T$ tel que $T \in [AB]$.
    \item Placer un point $N$ tel que $N \in [AB)$ et $N \notin [AB]$.
    \item Placer un point $C$ tel que $C \in [BA)$ et $C \notin [AB]$.
    \item Placer un point $H$ tel que $H \in [CA]$.
\end{compactitem}

\

\question[3] Sur votre feuille, réaliser les constructions demandées.
\begin{parts}
      \part[2] Tracer un cercle de centre $I$ et de diamètre 6 cms. Placer deux points J et K tels que le segment $[JK]$ soit un diamètre du cercle C, un point L sur ce même cercle C, un point M qui appartient au disque, et un point N qui n'appartient pas au disque.
      \part[0.5] \textbf{\textcolor{red}{[Justifier]}} Que représente le point I pour le segment $[JK]$ ? Coder la figure en conséquence.
      \part[0.5] \textbf{\textcolor{red}{[Justifier]}} Que vaut la longueur IL ? Coder la figure en conséquence.
  \end{parts}

\end{questions}

\section*{Exercice 3 - Cercle (3 points)}

\begin{questions}
  \question[2.5] Déterminer si les affirmations suivantes sont vraies ou fausses. Justifier la réponse en vous aidant de la figure 2.
  \begin{parts}
    \part[0.25] \textbf{\textcolor{red}{[Justifier]}} $B$ est à égale distance de $A$ et de $E$.
    \part[0.25] \textbf{\textcolor{red}{[Justifier]}} $G$ est à égale distance de $A$ et de $E$.
    \part[0.25] \textbf{\textcolor{red}{[Justifier]}} $G$ est le milieu de $[AE]$.
    \part[0.25] \textbf{\textcolor{red}{[Justifier]}} $D$ est le milieu de $[CE]$.
    \part[0.25] \textbf{\textcolor{red}{[Justifier]}} $F$ est le milieu de $[DF]$.
    \part[0.25] \textbf{\textcolor{red}{[Justifier]}} $C$ est le milieu de $[AE]$.
    \part[0.25] \textbf{\textcolor{red}{[Justifier]}} $[AE]$ est un diamètre du cercle.
    \part[0.25] \textbf{\textcolor{red}{[Justifier]}} $[AG]$ est un rayon du cercle.
    \part[0.25] \textbf{\textcolor{red}{[Justifier]}} $[GE]$ est une corde du cercle.
    \part[0.25] \textbf{\textcolor{red}{[Justifier]}} $A$ et $B$ appartiennent au disque 
  \end{parts}
\question[0.5] \textbf{\textcolor{red}{[Justifier]}} Comment facilement tracer la médiatrice du segment $[AE]$ ?
\end{questions}

\begin{figure}[H]
  \centering
  \includegraphics[width=0.5\linewidth]{figure_2.jpg}
  \captionsetup{labelformat=empty}
  \caption{\label{} Figure 2}
\end{figure} 

\section*{Exercice 4 - Trajet en avion (2 points)}

\begin{questions} 
  \question[2] \textbf{\textcolor{red}{[Justifier]}} On a représenté sur la carte le trajet d'un avion. De quelles villes l'avion est-il toujours à égale distance ? Justifier (les mesures du dessin ne sont pas forcément parfaitement précises mais proches de la réalité).
\end{questions}

  \begin{figure}[H]
    \centering
    \includegraphics[width=0.5\linewidth]{figure_3.jpg}
    \captionsetup{labelformat=empty}
    \caption{\label{} Figure 3}
  \end{figure}



\section*{Exercice 5 - Calculs - Rappels (3 points)}

\begin{questions}
  \question[2] Effectuer les calculs suivants en détaillant les étapes.
  \begin{parts}
    \part[0.5] \textbf{\textcolor{red}{[Justifier]}} $A = (1 + 9) \times (14 - 3)$
    \part[0.5] \textbf{\textcolor{red}{[Justifier]}} $B = 3 \times [10 - (2 + 3)] + 5$
    \part[0.5] \textbf{\textcolor{red}{[Justifier]}} $C = 25 \times 4,7 \times 0,5 \times 4 \times 2$
    \part[0.5] \textbf{\textcolor{red}{[Justifier]}} $D = 2,5 \times 5 \times 4 \times 177 \times 2$
    \end{parts}
  \question[1] Déterminer l'ordre de grandeur (et préciser à quel ordre de grandeur vous arrondissez) de :
  \begin{parts}
    \part[0.5] \textbf{\textcolor{red}{[Justifier]}} $J = (88,4 + 121,6) \times 9,8$
    \part[0.5] \textbf{\textcolor{red}{[Justifier]}} $K = 2,94 \times 5,27 + 5,081$
  \end{parts}
\end{questions}


\end{document}
